\documentclass{article}
\usepackage{/globals/de}
\begin{document}
\section{Lineare Gleichungssysteme} 
Ein lineares Gleichungssystem ist ein System aus mehreren Gleichungen mit mehreren unbekannten Variablen. Durch das Lösen des Gleichungssytems können alle Werte aller Variablen die die Gleichungen des System erfüllen gefunden werden.

\subsection{Lösungsverfahren}
Viele, einfacherere Gleichungssysteme können durch das wiederholte anwenden bestimmter Verfahren gelöst werden.
\begin{description}
  \item[Einsetzungsverfahrens] Eine der Gleichungen wird nach einer Variable aufgelöst. Dessen Wert wird in eine andere Gleichung eingesetzt. Somit kommt die Variable nach der Aufgelöst wurde nicht mehr in der anderen Gleichung vor; die Anzahl der Unbekannten dieser Gleichung wurde um eins verringert.
  \item[Additionsverfahrens] Zwei Gleichungen werden sich angeguckt und unabhängig voneinander so multipliziert, dass wenn sie zusammeaddiert oder subtrahiert werden, eine Variable wegfällt. 
\end{description}
Durch wiederholtes Anwenden dieser Verfahren können die Gleichungen Schritt für Schritt vereinfacht und die Anzahl der Unbekannten in ihnen verkleinert werden und der Zahlenwert der Variablen ermittelt und weiter eingesetzt werden. 

\subsection{Lösungen}
Beim Lösen können Zahlenwerte für alle Variablen als Ergebnis rauskommen; dann werden diese einfach am Ende doppelt unterschrichen aufgeschrieben.

Kommt es während des Lösens zu einem wiederspruch, \zb eine Zahl sei eine andere Zahl, so gibt es keine Kombination an Variablenwerten für die alle Gleichungen des Systems erfüllt sind; das System hat keine Lösung.

Kommt es zu einer Identität, der Erkentnis das eine Variable oder eine Zahl sich selbst ist, so gibt es unendlich viele Lösungen. Diese können bei Bedarf genauer ermittelt werden, indem die Werte jeder Variable als Formel in Abhängigkeit von einer unbekannten aufgeschrieben wird.  
\end{document}